% !TEX program = xelatex
% 刘皓 — AI应用工程师简历
% 编译: xelatex resume.tex

\documentclass[11pt,a4paper]{article}

% ── Packages ──
\usepackage[UTF8]{ctex}
\usepackage[margin=18mm, top=16mm, bottom=14mm]{geometry}
\usepackage{fontspec}
\usepackage{xcolor}
\usepackage{titlesec}
\usepackage{enumitem}
\usepackage{hyperref}
\usepackage{tabularx}

% ── Fonts ──
\setmainfont{Inter}[
  Path=./fonts/,
  Extension=.ttf,
  UprightFont=*-Regular,
  BoldFont=*-Bold,
  ItalicFont=*-Italic,
]
\setCJKmainfont{Source Han Sans SC}[
  Path=./fonts/,
  Extension=.otf,
  UprightFont=*-Regular,
  BoldFont=*-Bold,
]
\setmonofont{JetBrains Mono}[Path=./fonts/,Extension=.ttf,UprightFont=*-Regular]

% ── Colors ──
\definecolor{text}{HTML}{1a1a2e}
\definecolor{textsecondary}{HTML}{4b5563}
\definecolor{accent}{HTML}{2563eb}
\definecolor{border}{HTML}{e5e7eb}
\definecolor{tagbg}{HTML}{eff6ff}
\definecolor{tagtext}{HTML}{1d4ed8}

% ── Hyperlinks ──
\hypersetup{
  colorlinks=true,
  linkcolor=accent,
  urlcolor=textsecondary,
}

% ── Section styling ──
\titleformat{\section}
  {\Large\bfseries\color{text}}
  {}
  {0em}
  {}
  [\vspace{-0.5em}\rule{\textwidth}{0.4pt}\vspace{0.3em}]

\titlespacing{\section}{0pt}{12pt}{6pt}

% ── Compact lists ──
\setlist{nosep,leftmargin=*,itemsep=2pt,topsep=2pt}
\setlist[itemize]{label={\color{accent}\small $\blacktriangleright$}}

% ── No page number ──
\pagestyle{empty}

\begin{document}

% ═══════════════ HEADER ═══════════════
\begin{center}
  {\Huge\bfseries\color{text} 刘 皓} \\[4pt]
  {\large\color{accent} AI应用工程师 \ —— \ LLM应用开发 / AI Agent工程} \\[8pt]
  {\small\color{textsecondary}
    \href{mailto:1281741407@qq.com}{1281741407@qq.com} \quad
    \href{https://github.com/p0styyyc}{github.com/p0styyyc} \quad
    广州
  }
\end{center}

\vspace{4pt}

% ═══════════════ 技术能力 ═══════════════
\section{技术能力}

\newcommand{\skilltag}[1]{\colorbox{tagbg}{\textcolor{tagtext}{\texttt{\small #1}}}}

{\small
\textbf{LLM应用} \quad LangGraph · LangChain · Agent架构 · Tool Calling · Prompt Engineering · Multi-Agent协作 · RAG

\textbf{后端开发} \quad Python 3.11+ · FastAPI · WebSocket · RESTful API · Pydantic · 异步编程 · Docker SDK

\textbf{工程工具} \quad Git · Docker · Docker Compose · pytest · ruff · mypy · Linux

\textbf{前端（辅助）} \quad React 18 · TypeScript · Tailwind CSS
}

% ═══════════════ 项目经历 ═══════════════
\section{项目经历}

\noindent
{\large\bfseries Autonomous AI Software Development Agent} \hfill {\small\color{textsecondary}2026.06} \\
{\small\color{textsecondary} 个人作品集项目 · 70+ 文件 · 92 测试 · \href{https://github.com/p0styyyc/Autonomous-AI-Software-Development-Agent}{github.com/p0styyyc/Autonomous-AI-Software-Development-Agent}}

\vspace{4pt}

{\small 对标 GitHub Copilot Workspace / Devin 理念，从零构建完整多Agent自主编程系统。用户输入自然语言需求后，系统自动完成 \textbf{需求理解 → 任务规划 → 代码生成 → 沙箱执行 → 质量审查} 的完整闭环。}

\vspace{4pt}

\begin{itemize}[leftmargin=10pt]
  \item \textbf{Agent编排：}基于 LangGraph StateGraph 构建 Planner → Coder → Executor → Reviewer 四Agent流水线，TypedDict 单一状态源实现 Agent 间信息传递，条件路由实现 通过→继续 / 失败→重试(≤3次) / 超限→跳过 的自主决策循环

  \item \textbf{LLM工程化：}工厂模式适配 OpenAI / DeepSeek 双 Provider，不同 Agent 分配不同模型与温度参数（Planner: GPT-4o-mini 0.3 / Coder: GPT-4o 0 / Reviewer: GPT-4o-mini 0）。实现多格式 JSON 提取 + 字段补全 + 降级策略，解决 LLM 输出不可靠问题

  \item \textbf{Tool Calling体系：}实现 5 个定制 LangChain Tool（文件创建/读取/编辑/代码执行/测试运行），Registry 模式 + 角色最小权限分发。Docker 沙箱 5 层安全隔离：容器级隔离 · 网络禁用 · 只读根文件系统 · cap\_drop=ALL · 512MB/30s限制

  \item \textbf{实时通信：}基于 asyncio.Queue 的 Pub/Sub 事件总线，WebSocket 推送 13 种 Agent 事件（planning→coding→executing→reviewing→done），支持历史事件补发 + 心跳保活

  \item \textbf{前端可视化：}React 18 + TypeScript + Tailwind CSS，实现步骤进度条、文件树、代码高亮、最终报告卡片。Docker Compose 一键部署（Nginx 反向代理 + WebSocket 升级）
\end{itemize}

\vspace{2pt}
{\footnotesize\color{textsecondary} \textbf{技术栈} \quad Python · FastAPI · LangGraph · LangChain · React 18 · TypeScript · Docker · pytest · WebSocket · Pydantic}

% ═══════════════ 教育经历 ═══════════════
\section{教育经历}

\noindent
{\bfseries 广东药科大学} \ —— \ 电子信息工程 · 本科 \hfill {\small\color{textsecondary}2026届}

% ═══════════════ 个人优势 ═══════════════
\section{个人优势}

\begin{itemize}[leftmargin=10pt]
  \item \textbf{全链路AI工程能力：}独立完成 Agent架构设计 → LLM编排 → Tool Calling → 前后端开发 → Docker部署的完整闭环，非 Demo 级代码

  \item \textbf{工程化思维：}92 个单元/集成测试 · 全量类型注解 · Pydantic 校验 · 错误降级策略 · 安全防护机制 · ADR 架构决策记录

  \item \textbf{技术广度与深度：}从 LLM 应用到全栈开发到容器化部署，理解每一层的工作原理而非停留在调用层面
\end{itemize}

\end{document}
